\section{Uncertainty and Significance of a Skill Score}
\label{sec:uncertainty}

The metrics of Section~\ref{sec:metrics} are point estimates. A skill
score of, say, $-0.078$ reports that a model's average error is 7.8 percent
larger than persistence's on the evaluation sample, but it does not say
whether that gap is distinguishable from zero or is within sampling noise.
For a comparison to support a claim, the gap needs an interval and a
significance statement. This section formalizes both, implemented as
optional functions in \texttt{src/evaluate.py} that are deliberately not
wired into the mandatory comparison table.

\subsection{Loss differential}

Let $e_i^{\mathrm{m}} = |y_i - \hat y_i^{\mathrm{m}}|$ and
$e_i^{\mathrm{p}} = |y_i - \hat y_i^{\mathrm{p}}|$ be the absolute errors of
a candidate model and of persistence on row $i$. Their difference

\begin{equation}
\label{eq:loss-diff}
d_i = e_i^{\mathrm{p}} - e_i^{\mathrm{m}}
\end{equation}

is positive when the model beats persistence on row $i$. The classical
Diebold--Mariano test \cite{diebold1995comparing} assesses the null
hypothesis $\mathbb{E}[d_i]=0$ (equal predictive accuracy) using the mean
and a serial-correlation-robust variance of $\{d_i\}$. The skill score of
Equation~\ref{eq:skill} is a monotone reparameterization of the ratio of
mean absolute errors, so testing whether skill differs from zero is the
same question as testing whether the mean loss differential
(Equation~\ref{eq:loss-diff}) differs from zero.

\subsection{Why a row-level interval is wrong here}

The Diebold--Mariano test was designed for a single time series of
forecasts. This project's evaluation sample is a panel: each country
contributes several years, and a country's own years are correlated
(its emissions, and its errors, move together). Treating the $n$ rows as
independent would understate the sampling variance, because the effective
number of independent observations is closer to the number of countries
than to the number of rows. In addition, a few large emitters dominate the
absolute error (Section~\ref{sec:metrics}), so the influence of individual
countries is uneven. Both features call for resampling at the level of the
country, not the row.

\subsection{Country-clustered bootstrap interval}

Let $\mathcal{C}$ be the set of countries in the evaluation sample, with
$|\mathcal{C}| = C$. A cluster bootstrap \cite{cameron2008bootstrap},
\cite{efron1993bootstrap} draws $C$ countries with replacement, pools all
rows of the drawn countries, and recomputes the statistic. For bootstrap
replicate $b = 1,\dots,B$, let $\mathcal{C}^{(b)}$ be the multiset of drawn
countries and $R^{(b)}$ the pooled rows; the resampled skill is

\begin{equation}
\label{eq:boot-skill}
s^{(b)} = 1 - \frac{\sum_{i \in R^{(b)}} e_i^{\mathrm{m}}}
                   {\sum_{i \in R^{(b)}} e_i^{\mathrm{p}}}
\end{equation}

Resampling whole countries preserves the within-country dependence that a
row-level resample would destroy. A $100(1-\alpha)$ percent confidence
interval for the skill is the empirical interval of the replicates,

\begin{equation}
\label{eq:boot-ci}
\left[\; q_{\alpha/2}\!\left(\{s^{(b)}\}\right),\;
        q_{1-\alpha/2}\!\left(\{s^{(b)}\}\right) \;\right]
\end{equation}

where $q_\tau$ is the $\tau$ empirical quantile. A two-sided bootstrap
$p$-value for the null hypothesis that the model ties persistence
($s = 0$) is

\begin{equation}
\label{eq:boot-p}
p = 2 \min\!\left(
  \frac{1}{B}\sum_{b=1}^{B}\mathbf{1}\{s^{(b)} \le 0\},\;
  \frac{1}{B}\sum_{b=1}^{B}\mathbf{1}\{s^{(b)} \ge 0\}
\right),
\end{equation}

clipped to $[0,1]$. This is a bootstrap tail probability, not an exact
analytic test, and should be reported as such.

\subsection{Worked reading}

On this project's validation years (targets 2015--2018, 153 countries),
the linear-trend baseline has point skill $-0.078$ against persistence, but
the country-clustered 95 percent interval (Equation~\ref{eq:boot-ci}) spans
roughly $[-0.34, +0.20]$ with a bootstrap $p$-value near $0.7$. The honest
conclusion is therefore not ``trend loses to persistence'' but ``trend and
persistence are statistically indistinguishable on this sample'': the point
gap is well within sampling noise once the country-level dependence is
accounted for. Reporting the interval, not only the point estimate, is what
keeps the evaluation honest at the level of statistical significance, which
is the same discipline this project already applies at the level of the
baseline comparison itself.
